\exercisegroup

\begin{enumerate}
\def\labelenumi{\arabic{enumi})}
\tightlist
\item
  For every \(\theta \in \mathbf{R}\) , let \(C_{\theta}\) be the segment of the plane with endpoints \((0,0)\) and \((\cos(2\pi\theta), \sin(2\pi\theta))\) ; set
\end{enumerate}

\[
\mathrm{C} = \bigcup_ {\theta \in \mathbf {Q} \cap [ 1 / 6, 1 / 4 ]} \mathrm{C} _ {\theta}.
\]

\begin{enumerate}
\def\labelenumi{\alph{enumi})}
\item
  The space C is contractible at 0, but at no other point.
\item
  Let D be the union of the translates \((0, n) + \mathrm{C}\) of C, for \(n \in Z\) . Show that the space D is homotopy equivalent to a point but is not contractible at any of its points.
\end{enumerate}

\begin{enumerate}
\def\labelenumi{\arabic{enumi})}
\setcounter{enumi}{1}
\tightlist
\item
  Let B be a subset of \(\mathbf{R}^{n}\) , let \(a\) be a point interior to B; denote by S the boundary of B. Suppose that B is star-shaped at \(a\) and that every half-line with origin \(a\) meets S in a unique point. Let \(f\colon\mathbf{R}_{+}^{*}\times\mathrm{S}\to\mathbf{R}^{n}-\{a\}\) be the map defined by \(f(t,y)=ty+(1-t)a\) . Prove that \(f\) is a homeomorphism such that \(f([0,1]\times\mathrm{S})=\mathrm{B}-\{a\}\) and \(f(\{1\}\times\mathrm{S})=\mathrm{S}\) .
\end{enumerate}

This result applies in particular when B is a convex, closed, and bounded subset of \(R^{n}\) and a is a point adherent to B (cf. I, p. 123, prop. 2 and EVT, II, p. 15, prop. 16).

\begin{enumerate}
\def\labelenumi{\arabic{enumi})}
\setcounter{enumi}{2}
\item
  Let X be the topological space obtained from the space R by contracting the subset Q (III, p. 252, def. 10). The canonical surjection \(p: R \rightarrow X\) is neither open nor closed.
\item
  Let X be the interval \([-1,1]\) of R and let A = X - \{0\}. The canonical map of X onto X/A is a homotopy equivalence but there does not exist a homotopy σ with origin \(\mathrm{Id}_{X}\) satisfying the hypotheses of prop. 10 of III, p. 252.
\item
  Let A be the closure in \(\mathbf{R}\) of the set of real numbers of the form \(1/n\) , where \(n\) runs through the set of integers \(\geqslant 1\) . Let \(i\) be the canonical injection of A into \(\mathbf{R}\) . The canonical map of \(\mathrm{Côn}(i)\) onto \(\mathbf{R}/\mathrm{A}\) is not a homotopy equivalence.
\item
  Let X be a metrizable topological space of countable type. We say that X is an absolute retract\(^{(1)}\) if for every metrizable topological space of countable type Y and every closed subset B of Y, every continuous map \(f: B \to X\) extends to a continuous map from Y into X.
\end{enumerate}

\begin{enumerate}
\def\labelenumi{\alph{enumi})}
\item
  Let X be a metrizable topological space of countable type. For X to be an absolute retract, it is necessary and sufficient that, for every metrizable space of countable type Y and every continuous, injective, and closed map f from X into Y, there exists a continuous map \(g: Y \to X\) such that \(g \circ f = Id_{X}\) .
\item
  Show that the spaces I and R are absolute retracts.
\item
  Show that the product space of a countable family of absolute retracts is an absolute retract.
\item
  Let X be an absolute retract, let A be a subspace of X which is a retract of X. Show that A is an absolute retract.
\end{enumerate}

\begin{enumerate}
\def\labelenumi{\arabic{enumi})}
\setcounter{enumi}{6}
\tightlist
\item
  Let X be a metrizable topological space of countable type. We say that X is an absolute neighborhood retract if for every metrizable topological space of countable type Y and every closed subset B of Y, every continuous map \(f: B \rightarrow X\) extends to a continuous map from a neighborhood of B into X.
\end{enumerate}

\begin{enumerate}
\def\labelenumi{\alph{enumi})}
\item
  Let X be a metrizable topological space of countable type. For X to be an absolute neighborhood retract, it is necessary and sufficient that for every metrizable space of countable type Y and every continuous, injective, and closed map f from X into Y, there exists a neighborhood U of \(f(\mathrm{X})\) and a continuous map \(g: U \to X\) such that \(g \circ f = Id_{X}\) .
\item
  Let X be an absolute neighborhood retract and let A be a subset of X. If A has a neighborhood U of which it is a retract, then A is an absolute neighborhood retract. In particular, every open subset of X and every retract of X is an absolute neighborhood retract.
\item
  Show that the product space of a finite family of absolute neighborhood retracts is an absolute neighborhood retract.
\end{enumerate}

d) Show that the sphere \(\mathbf{S}_n\) is an absolute neighborhood retract.

\begin{enumerate}
\def\labelenumi{\alph{enumi})}
\setcounter{enumi}{4}
\item
  Let X be a metrizable topological space of countable type, let \(A_{1}\) and \(A_{2}\) be closed subsets of X such that \(X = A_{1} \cup A_{2}\) . We assume that \(A_{1} \cap A_{2}\) is an absolute neighborhood retract. For X to be an absolute neighborhood retract, it is necessary and sufficient that the same be true of \(A_{1}\) and of \(A_{2}\) .
\item
  Let X be a metrizable topological space of countable type. If X is the union of two open subsets which are absolute neighborhood retracts, then X is an absolute neighborhood retract.
\item
  Let X be a metrizable topological space of countable type such that every point has a neighborhood which is an absolute neighborhood retract. Show that X is an absolute neighborhood retract.
\end{enumerate}

\begin{enumerate}
\def\labelenumi{\arabic{enumi})}
\setcounter{enumi}{7}
\tightlist
\item
  Let X be a metrizable topological space; let d be a metric defining its topology. Let \(\mathcal{C}_{b}(\mathrm{X};\mathbf{R})\) be the vector space of real-valued functions on X that are continuous and bounded; we equip it with the norm defined by \(\|f\| = \sup_{x \in X} |f(x)|\) for \(f \in \mathcal{C}_{b}(\mathrm{X};\mathbf{R})\) .
\end{enumerate}

\begin{enumerate}
\def\labelenumi{\alph{enumi})}
\item
  One defines a map \(j\colon X \to \mathcal{C}_{b}(X; \mathbf{R})\) by setting \(j(x)\colon y \mapsto d(x, y)/(1 + d(x, y))\) . Prove that the map j induces a homeomorphism from X onto its image.
\item
  Let V be the intersection of all convex subsets of \(\mathcal{C}_{b}(\mathrm{X};\mathbf{R})\) that contain \(j(\mathrm{X})\) . Prove that \(j(\mathrm{X})\) is closed in V.
\item
  Let X be a metrizable space of countable type. Prove that there exists a convex subset Z of \(I^{N}\) such that X is homeomorphic to a closed subset of Z.
\end{enumerate}

\begin{enumerate}
\def\labelenumi{\arabic{enumi})}
\setcounter{enumi}{8}
\tightlist
\item
  If X and Y are topological spaces and U is a set of subsets of X, we say that maps \(f, g: Y \to X\) are U-close if for every point \(y \in Y\) , there exists \(U \in U\) such that \(f(y)\) and \(g(y)\) belong to U; we say that a homotopy \(\sigma: Y \times I \to X\) is U-small if for every point \(y \in Y\) , there exists \(U \in U\) such that \(\sigma(y, t) \in U\) for all \(t \in I\) .
\end{enumerate}

Let X be a closed subspace of a convex subset Z of \(I^{N}\) (cf. Exercise 8); we assume that X is an absolute neighborhood retract.

\begin{enumerate}
\def\labelenumi{\alph{enumi})}
\item
  Let U be a cover of X by open subsets of \(I^{N}\) . Prove that there exists a family V of open and convex subsets of \(I^{N}\) such that, for every open set \(v \in V\) , there exists an open set \(u \in U\) such that \(v \cap X \subset u\) , and a continuous retraction r of the canonical injection of X into the union V of the family V.
\item
  Let Y be a topological space and let B be a closed subspace of Y. Let \(f, g: \mathrm{Y} \to \mathrm{X}\) be continuous maps that are \(\mathcal{V}\)-close, and let \(\sigma: \mathrm{B} \times \mathbf{I} \to \mathrm{X}\) be a homotopy \(\mathcal{V}\)-small whose origin is equal to \(f|_{\mathrm{B}}\) .
\end{enumerate}

Prove that there exists a neighborhood W of B in Y and a continuous map h from the subspace \((\mathbf{Y} \times \{0,1\}) \cup (\mathbf{W} \times \mathbf{I})\) of \(Y \times I\) into X such that \(h(y,0) = f(y)\) and \(h(y,1) = g(y)\) for all \(y \in Y\) , and \(h(y,t) = \sigma(y,t)\) for all \(y \in B\) and every \(t \in I\) . Prove that there exists a neighborhood \(W'\) of B contained in W such that \(h \mid W' \times I\) is a \(\mathcal{V}\)-small homotopy.

\begin{enumerate}
\def\labelenumi{\alph{enumi})}
\setcounter{enumi}{2}
\tightlist
\item
  Let \(u \colon \mathrm{Y} \to \mathbf{I}\) be a continuous map such that \(u(y) = 1\) for \(y \in \mathrm{B}\) and \(u(y) = 0\) in a neighborhood of \(\mathrm{Y} - \mathrm{W}\) ; set
\end{enumerate}

\[
h ^ {\prime} (y, t) = \left\{ \begin{array}{l l} (1 - u (y)) ((1 - t) f (y) + t g (y)) + u (y) h (y, t) & \text { if } y \in \mathrm{V} \\ (1 - t) f (y) + t g (y) & \text { otherwise. } \end{array} \right.
\]

Prove that \(h'\) is continuous and that its image is contained in V.

\(d)\) Deduce that there exists a homotopy \(\widetilde{\sigma}\colon\mathbf{Y}\times\mathbf{I}\to\mathbf{X}\) with source \(f\) , with endpoint \(g\) , which is \(\mathcal{U}\) -small.

\begin{enumerate}
\def\labelenumi{\arabic{enumi})}
\setcounter{enumi}{9}
\tightlist
\item
  Let X be a topological space. Suppose that there exists an open cover U of X such that for every topological space Y and every closed subspace B of Y, two maps \(f_{0}, f_{1}\) from Y to X which are U-close are homotopic if there exists a U-small homotopy connecting \(f_{0} \mid B\) to \(f_{1} \mid B\) .
\end{enumerate}

\begin{enumerate}
\def\labelenumi{\alph{enumi})}
\item
  Let a be a point of X and let U be an element of U such that \(a \in U\) . Prove that there exists a homotopy \(\sigma: U \times I \to X\) such that \(\sigma(x,0) = a\) , \(\sigma(x,1) = x\) and \(\sigma(a,t) = a\) for all \(x \in X\) and every \(t \in I\) .
\item
  Let V be an open neighborhood of a such that \(\sigma(\mathrm{V} \times \mathbf{I}) \subset \mathrm{U}\) . Prove that V is an absolute neighborhood retract.
\item
  Prove that X is an absolute neighborhood retract.
\end{enumerate}

\begin{enumerate}
\def\labelenumi{\arabic{enumi})}
\setcounter{enumi}{10}
\tightlist
\item
  Let X be an absolute neighborhood retract.
\end{enumerate}

\begin{enumerate}
\def\labelenumi{\alph{enumi})}
\item
  Prove that X is "semi-locally contractible", that is, every point x of X has a neighborhood V such that the canonical injection of (V, x) into X is strictly homotopic to the constant map with image x. (Use Exercise 9.)
\item
  Let A be a compact topological space. Prove that the path-connected components of the space \(\mathcal{C}_{c}(\mathrm{A};\mathrm{X})\) are open.
\item
  Prove that X is locally path-connected.
\end{enumerate}

\begin{enumerate}
\def\labelenumi{\arabic{enumi})}
\setcounter{enumi}{11}
\tightlist
\item
  Let X and Y be topological spaces; one says that the homotopy extension property holds for continuous maps with target X and source Y if for every continuous map F: \(Y \rightarrow X\) , every closed subspace B of Y, and every homotopy \(\sigma\colon\mathrm{B}\times\mathbf{I}\to\mathrm{X}\) with origin F \textbar{} B, there exists a homotopy \(\widetilde{\sigma}\colon\mathrm{Y}\times\mathbf{I}\to\mathrm{X}\) with origin F which extends \(\sigma\) .
\end{enumerate}

\begin{enumerate}
\def\labelenumi{\alph{enumi})}
\item
  Let X be a topological space which is an absolute neighborhood retract. Prove that the homotopy extension property holds for continuous maps with target X and source a metrizable space of countable type.
\item
  Prove the equivalence of the following two properties:
\end{enumerate}

\begin{enumerate}
\def\labelenumi{(\roman{enumi})}
\item
  The space X is an absolute neighborhood retract;
\item
  Every point of X has a neighborhood V such that for every metrizable topological space Y of countable type, every closed subspace B of Y, and every continuous map from B to V extends to a continuous map from Y to X.
\end{enumerate}

\begin{enumerate}
\def\labelenumi{\alph{enumi})}
\setcounter{enumi}{2}
\item
  Suppose that X is semi-locally contractible (cf. Exercise 9) and that the homotopy extension property holds for continuous maps with target X. Prove that X is an absolute neighborhood retract. \(^{(2)}\)
\item
  Prove that the space X = Q is not an absolute neighborhood retract but that the homotopy extension property holds for continuous maps with target X.
\end{enumerate}

\begin{enumerate}
\def\labelenumi{\arabic{enumi})}
\setcounter{enumi}{12}
\item
  Let X be a topological space which is an absolute neighborhood retract. Let A be a subset of X such that the canonical injection of A into X is a homotopy equivalence and a cofibration. There exists a contraction of X onto A.
\item
  Let X be the topological space I and set A = \{0,1\}. Prove that there exists a homotopy connecting the identity map of X to a constant map but that the canonical map X → X/A is not a homotopy equivalence.
\item
  Let X be the subspace of \([0,1]\times[0,1]\) consisting of the pairs \((x,y)\) such that \(x\in\{0,1\}\) , or \(y\in\{0,1\}\) , or \(x\neq0\) and \(1/x\in N\) . Let A = \(\{0\}\times[0,1]\) . Show that the canonical map from X onto X/A is not a homotopy equivalence, although A is contractible.
\item
  Let X and Y be topological spaces, and let \(f: X \to Y\) be a continuous map. Let \(\varphi: C(X) \to \text{Côn}(f)\) be the canonical map. Show that the spaces \(Y/f(X)\) and \(\text{Côn}(f)/\varphi(C(X))\) are homeomorphic.
\item
  Let X be a topological space and let A be a subspace of X. Denote by \(i\colon A \to X\) and \(j\colon \operatorname{Cyl}(i) \to X \times I\) the canonical injections.
\end{enumerate}

\begin{enumerate}
\def\labelenumi{\alph{enumi})}
\item
  Suppose that X = R and \(A = ]0, 1[\); show that the map j is not strict.
\item
  Suppose that every point of X has a countable base of closed neighborhoods. For the map j to be strict, it is necessary and sufficient that A be closed.
\item
  Suppose that A is the Alexandroff half-line (TG, IV, p. 49, exerc. 12) and that X is its Alexandroff compactification. Show that the map j is strict although A is not closed in X.
\end{enumerate}

\begin{enumerate}
\def\labelenumi{\arabic{enumi})}
\setcounter{enumi}{17}
\tightlist
\item
  A topological space X is said to be locally equiconnected if the diagonal map from X into X × X is a cofibration.
\end{enumerate}

\begin{enumerate}
\def\labelenumi{\alph{enumi})}
\tightlist
\item
  Show that the product of a finite family of locally equiconnected spaces is locally equiconnected.
\end{enumerate}

Let X be a locally equiconnected space.

\begin{enumerate}
\def\labelenumi{\alph{enumi})}
\setcounter{enumi}{1}
\item
  Show that there exists a covering of X by a sequence \((\mathrm{U}_{n})_{n\in\mathbf{N}}\) of open subsets of X such that, for every n, the injection of \(U_{n}\) into X is homotopic to a constant map.
\item
  Show that X is Hausdorff.
\item
  Let \(u: X \to I\) be a continuous map and let U be the set of points \(x \in X\) such that \(u(x) > 0\) . Show that the path-connected components of U are open. In particular, the path-connected components of X are open.
\item
  Let A be a closed subspace of X which is a retract of X. Show that A is locally equiconnected and that the pair (X, A) has the homotopy extension property. ("Theorem of Dyer and Eilenberg".)
\item
  Let A be a subspace of X such that the pair (X, A) has the homotopy extension property. Show that A is locally equiconnected.
\end{enumerate}

\begin{enumerate}
\def\labelenumi{\arabic{enumi})}
\setcounter{enumi}{18}
\tightlist
\item
  Let B be a topological space, let \(u \colon \mathrm{X} \to \mathbf{I}\) be a continuous map and let \(\mathrm{A} = u^{-1}(0)\) .
\end{enumerate}

\begin{enumerate}
\def\labelenumi{\alph{enumi})}
\item
  Show that the map p from \(X \times I\) into itself given by \(p(x,t) = (x,tu(x))\) is closed.
\item
  Let Y be a topological space and let \(\sigma\colon\mathrm{X}\times\mathbf{I}\to\mathrm{Y}\) be a homotopy fixed on A. Let \(\tau\colon\mathrm{X}\times\mathbf{I}\to\mathrm{Y}\) be the map given by \(\tau(x,t)=\sigma(x,\inf(1,t/u(x)))\) if \(u(x)>0\) and \(\tau(x,t)=\sigma(x,1)\) otherwise, prove that \(\tau\) is continuous.
\item
  Let C be a subspace of X; suppose that the pair (X, C) has the homotopy extension property. Let \(f \colon \mathrm{X} \to \mathrm{Y}\) be a continuous map and let \(\sigma \colon \mathrm{C} \times \mathbf{I} \to \mathrm{Y}\) be a homotopy with origin \(f \mid \mathrm{C}\) which is fixed on \(\mathrm{A} \cap \mathrm{C}\) . Prove that there exists a homotopy \(\tau \colon \mathrm{X} \times \mathbf{I} \to \mathrm{Y}\) with source \(f\) which extends \(\sigma\) and which is fixed on \(\mathrm{A}\) .
\end{enumerate}

\begin{enumerate}
\def\labelenumi{\arabic{enumi})}
\setcounter{enumi}{19}
\tightlist
\item
  Let X be a topological space and let A, B, and C be closed subspaces of X such that \(A = B \cap C\) .
\end{enumerate}

\begin{enumerate}
\def\labelenumi{\alph{enumi})}
\item
  If the pairs (X, B), (X, C), and (X, A) have the homotopy extension property, prove that the pair (X, B ∪ C) also has it.
\item
  Assume that the pairs (X, A) and (X, B ∪ C) have the homotopy extension property. Prove that the pairs (X, B) and (X, C) also have it.
\item
  Assume that X is a normal space and that \(A \subset \mathring{B} \cup \mathring{C}\) . If the pairs (X, B) and (X, C) have the homotopy extension property, prove that the pair (X, A) also has it.
\end{enumerate}

\begin{enumerate}
\def\labelenumi{\arabic{enumi})}
\setcounter{enumi}{20}
\tightlist
\item
  Let X be a paracompact topological space, and let \(\mathcal{A}\) be a set of closed subsets of X such that \(B \cap C \in \mathcal{A}\) for every pair (B, C) of elements of \(\mathcal{A}\). Let \(A = \bigcup_{C \in \mathcal{A}} C\) ; we assume that, for every point x of X, there exists a neighborhood V of x and a finite subset \(\mathcal{A}_x\) of \(\mathcal{A}\) such that \(V \cap A = V \cap \bigcup_{C \in \mathcal{A}_x} C\) . a) Prove that A is a closed subset of X.
\end{enumerate}

\begin{enumerate}
\def\labelenumi{\alph{enumi})}
\setcounter{enumi}{1}
\tightlist
\item
  We assume that, for all \(B \in A\) , the pair (X, B) has the homotopy extension property. Prove that the pair (X, A) also has this property.
\end{enumerate}

\begin{enumerate}
\def\labelenumi{\arabic{enumi})}
\setcounter{enumi}{21}
\tightlist
\item
  Let X and Y be topological spaces, let A be a subspace of X, and let B be a subspace of Y.
\end{enumerate}

\begin{enumerate}
\def\labelenumi{\alph{enumi})}
\item
  Assume that the pairs (X,A) and (Y,B) are homotopy equivalent. If the pair (X,A) has the homotopy extension property, prove that the pair (Y,B) also has it.
\item
  Conversely, suppose that the pairs (X, A) and (Y, B) have the homotopy extension property. Let \(f \colon X \to Y\) be a homotopy equivalence which induces, by restriction to subspaces, a homeomorphism from A onto B. Prove that \(f\) is a homotopy equivalence of the pair (X, A) to the pair (Y, B).
\end{enumerate}

\begin{enumerate}
\def\labelenumi{\arabic{enumi})}
\setcounter{enumi}{22}
\tightlist
\item
  Let X and Y be topological spaces. A homotopy \(\sigma\colon\mathrm{X}\times\mathbf{I}\to\mathrm{Y}\) is said to be stationary if there exists a real number \(\delta>0\) such that \(\sigma(x,t)=\sigma(x,0)\) for all \(x\in\mathrm{X}\) and every \(t\in[0,\delta]\) . Let A be a closed subspace of X; we denote by \(j\) the canonical injection of A into X. The pair (X,A) is said to have the homotopy extension property for stationary homotopies if, for every topological space Y, every continuous map \(f\colon\mathrm{X}\to\mathrm{Y}\) and every stationary homotopy \(\sigma\colon\mathrm{A}\times\mathbf{I}\to\mathrm{Y}\) whose origin is equal to \(f\mid\mathrm{A}\) , there exists a homotopy \(\widetilde{\sigma}\colon\mathrm{X}\times\mathbf{I}\to\mathrm{Y}\) with source \(f\) which extends \(\sigma\) .
\end{enumerate}

\begin{enumerate}
\def\labelenumi{\alph{enumi})}
\tightlist
\item
  Prove that the following properties are equivalent:
\end{enumerate}

\begin{enumerate}
\def\labelenumi{(\roman{enumi})}
\item
  The pair (X, A) has the stationary homotopy extension property.
\item
  (X, A) and \((\mathrm{Cyl}(j), \alpha_{j}(\mathrm{A} \times \{0\}))\) are homotopy equivalent;
\item
  There exists a continuous map \(\varphi\colon X\to\mathbf{I}\) such that \(\varphi(a)=0\) for all \(a\in A\) and a homotopy \(\sigma\colon X\times\mathbf{I}\to X\) such that \(\sigma(x,0)=x\) for \(x\in X\) , \(\sigma(a,t)=a\) for \((a,t)\in A\times\mathbf{I}\) , and \(\sigma(x,1)\in A\) if \(\varphi(x)<1\) .
\item
  There exists a continuous map \(\varphi\colon\mathrm{X}\to\mathbf{I}\), a subspace V of X, and a homotopy \(\sigma\colon\mathrm{V}\times\mathbf{I}\to\mathrm{X}\) such that \(\varphi(x)=1\) if \(x\notin\mathrm{V}\) , \(\varphi(a)=0\) for \(a\in\mathrm{A}\) , \(\sigma(x,0)=x\) for \(x\in\mathrm{V}\) , \(\sigma(a,t)=a\) for \((a,t)\in\mathrm{A}\times\mathbf{I}\) , and \(\sigma(x,1)\in\mathrm{A}\) if \(\varphi(x)<1\) .
\end{enumerate}

\begin{enumerate}
\def\labelenumi{\alph{enumi})}
\setcounter{enumi}{1}
\tightlist
\item
  We assume that the pair (X, A) has the property of extension of stationary homotopies and that there exists a continuous function \(u \colon X \to I\) such that \(u^{-1}(0) = A\) . Show that the pair (X, A) has the homotopy extension property.
\end{enumerate}

\begin{enumerate}
\def\labelenumi{\arabic{enumi})}
\setcounter{enumi}{23}
\tightlist
\item
  Let M be an uncountable set, let X be the space \(I^{M}\) and let A be the subspace \(\{0\}^{M}\) of X.
\end{enumerate}

\begin{enumerate}
\def\labelenumi{\alph{enumi})}
\item
  Prove that the pair (X, A) has the homotopy extension property for stationary homotopies (exercise 23).
\item
  Show that the pair (X, A) does not have the homotopy extension property. (Prove that the set A is not the intersection of a countable family of open sets.)
\item
  Prove that the pair \((\mathbf{X} \times \mathbf{I}, \mathbf{X} \times \{0\} \cup \mathbf{A} \times \mathbf{I})\) does not have the property of extension of stationary homotopies.
\end{enumerate}

\begin{enumerate}
\def\labelenumi{\arabic{enumi})}
\setcounter{enumi}{24}
\tightlist
\item
  \begin{enumerate}
  \def\labelenumii{\alph{enumii})}
  \tightlist
  \item
    Let X and Y be topological spaces, let A be a closed subspace of X, and let B be a closed subspace of Y. Let \(f \colon \mathrm{X} \to \mathrm{Y}\) be a continuous map; we assume that \(f\) induces, by passing to subspaces, a homeomorphism from A onto B. If \(f\) has a left inverse up to homotopy and if (Y, B) has the homotopy extension property for stationary homotopies (exerc. 23), then so does (X, A).
  \end{enumerate}
\end{enumerate}

\begin{enumerate}
\def\labelenumi{\alph{enumi})}
\setcounter{enumi}{1}
\item
  If the pairs (X, A) and (Y, B) are homotopy equivalent, then (X, A) has the stationary homotopy extension property if and only if (Y, B) does as well.
\item
  Let C be the set of points in the plane of the form \((t/n,t)\) , for \(t \in I\) and \(n \in N^{*}\) ; let X be its closure and let \(A = \{0\} \times I\) . Show that the pair \((X,A)\) does not have the property of extension of stationary homotopies.
\end{enumerate}

\begin{enumerate}
\def\labelenumi{\arabic{enumi})}
\setcounter{enumi}{25}
\item
  We equip the set \(X = \{0,1\}\) with the topology for which the open sets are \(\varnothing\) , \(\{0\}\) and X. We set \(A = \{0\}\) . Show that the pair \((X,A)\) has the homotopy extension property but the pair \((X \times X,(X \times A) \cup (A \times X))\) does not possess it.
\item
  Let X be a topological space and let A be a subspace (not necessarily closed) of X; denote by C the subspace \((\mathrm{X} \times \{0\}) \cup (\mathrm{A} \times \mathbf{I})\) of \(X \times I\) . We assume that C is a retract of \(X \times I\) . Let U be a subset of C such that \(\mathrm{U} \cap (\mathrm{X} \times \{0\})\) and \(\mathrm{U} \cap (\mathrm{A} \times \mathbf{I})\) are open in \(X \times \{0\}\) and \(A \times I\) respectively. a) Let \(V_{0}\) be the set of \(x \in X\) such that \((x, 0) \in \mathrm{U}\) ; for \(n \geqslant 1\) , let \(V_{n}\) be the union of the open sets W of X such that \((\mathrm{W} \cap \mathrm{A}) \times [0, 1/n]\) is contained in U.
\end{enumerate}

Prove that \(A \cap V_{0} = \bigcup_{n \geqslant 1} A \cap V_{n}\) ; also prove that every open set W of X such that \(W \cap A \subset V_{n}\) is contained in \(V_{n}\) .

\begin{enumerate}
\def\labelenumi{\alph{enumi})}
\setcounter{enumi}{1}
\item
  Prove that \(\mathrm{V} \subset \bigcup_{n \geqslant 1} \mathrm{V}_n\) .
\item
  Prove that U is open in C.
\item
  Prove that the pair (X, A) has the homotopy extension property.
\end{enumerate}

\begin{enumerate}
\def\labelenumi{\arabic{enumi})}
\setcounter{enumi}{27}
\item
  Let \(X\) be a topological space and let \(A\) be a subspace of X such that the pair \((X, A)\) has the homotopy extension property. Prove that the pair \((X, \overline{A})\) has the same property.
\item
  Let B be a topological space. A covering \((\mathrm{V}_{i})_{i\in\mathrm{I}}\) of B is numerable if there exists a continuous partition of unity \((g_{j})_{j\in\mathrm{J}}\) on B, locally finite, such that for every \(j\in J\) , \(g_{j}^{-1}(]0,1])\) is contained in one of the \(V_{i}\) .
\end{enumerate}

\begin{enumerate}
\def\labelenumi{\alph{enumi})}
\item
  Prove that B is paracompact if and only if every open cover of B is numerable.
\item
  Prove that B is normal if and only if every locally finite open cover of B is numerable.
\item
  Let \((\mathrm{V}_{i})\) be a numerable covering of B, let \(B'\) be a topological space and let \(f: B' \to B\) be a continuous map. Show that the covering \((f^{-1}(\mathrm{V}_{i}))\) of \(B'\) is numerable.
\end{enumerate}

\begin{enumerate}
\def\labelenumi{\arabic{enumi})}
\setcounter{enumi}{29}
\item
  Let B be a paracompact topological space, let F be a contractible topological space, and let E be a locally trivial fiber bundle over B with fiber F. Prove that the sheaf of germs of continuous sections of E is soft.
\item
  Let B be a topological space. Let A and V be subspaces of B; one says that V is a halo of A if there exists a continuous map \(\varphi\colon\mathbf{B}\to\mathbf{I}\) which is equal to 1 at every point of A and to 0 at every point of CV. A B-space (E,p) is said to have the section extension property if, for every subspace A of B, every halo V of A, and every section s of E|V, there exists a section s' of E such that s' \textbar{} A = s \textbar{} A.
\end{enumerate}

\begin{enumerate}
\def\labelenumi{\alph{enumi})}
\item
  Let \((\mathrm{E},p)\) and \((\mathrm{E}^{\prime},p^{\prime})\) be B-spaces; suppose that there exist morphisms of B-spaces, \(f: E \to E'\) and \(g: E' \to E\) , and a homotopy \(\sigma: E \times I \to E\) with source \(Id_{E}\) and endpoint \(g \circ f\) such that \(p(\sigma(x,t)) = p(x)\) for all \((x,t) \in E \times I\) . If the B-space \(E'\) has the section extension property, the same holds for the B-space E.
\item
  Let E be a B-space and let A be a subspace of B which is a retract of B. Suppose that the induced A-space \(E_{A}\) has the section extension property; prove that the B-space E has the same property.
\item
  Let E be a B-space which possesses the section extension property. Let \(f \colon \mathrm{B} \to \mathbf{I}\) be a continuous map and let \(V = f^{-1}(]0,1])\) . Prove that the V-space \(E_V\) has the section extension property. (Let A be a subspace of V and let W be a halo of A in V; construct sequences \((A_n)\) and \((V_n)\) of subspaces of B satisfying the following properties: for every \(n\) , \(A_n\) is contained in W, \(V_n\) is a halo of \(A_n\) , and the intersection of the \(A_n\) is equal to A.)
\end{enumerate}

\begin{enumerate}
\def\labelenumi{\arabic{enumi})}
\setcounter{enumi}{31}
\tightlist
\item
  Let B be a topological space and let \((\mathrm{E},p)\) be a B-space. Let \((\mathrm{V}_{i})_{i\in\mathrm{I}}\) be a numerable covering of B such that, for every \(i\in I\) , the \(V_{i}\)-space \(E_{V_{i}}\) has the section extension property.
\end{enumerate}

\begin{enumerate}
\def\labelenumi{\alph{enumi})}
\item
  We assume that there exists a continuous, locally finite partition of unity \((f_{i})\) on B such that \(V_{i} = f_{i}^{-1}(]0,1])\) for all i. Let \(g: B \to [0,1]\) be a continuous map. We set \(A = g^{-1}(1)\) and \(W = g^{-1}(]0,1])\) . Let s be a section of E over W. Prove that there exists a section of E that coincides with s on A. (Consider a maximal element of the set of pairs \((J,t)\) where J is a subset of I and t is a section of E over \(W \cup \bigcup_{i \in J} V_{i}\) which coincides with s on A, equipped with the order relation given by \((J,t) \prec (J',t')\) if \(J \subset J'\) and \(t'\) extends t.)
\item
  Prove that the B-space E has the section extension property.
\end{enumerate}

